% Job posting: https://ca.linkedin.com/jobs/view/audio-software-engineer-ea-sports%E2%84%A2-ufc-at-ea-sports-4452297229
\documentclass[11pt]{article}
\usepackage[left=1in,top=1in,right=1in,bottom=1in]{geometry}
\usepackage[hidelinks]{hyperref}
\usepackage{setspace}
\usepackage[T1]{fontenc}
\usepackage{newtxtext,newtxmath}
\ifdefined\pdfgentounicode
  \input{glyphtounicode}
  \pdfgentounicode=1
\fi
\setstretch{1.1}
\pagenumbering{gobble}
\begin{document}
\pagestyle{empty}

\begin{flushright}
Nishal Stanislaus Silva, PhD \\
Toronto, ON, Canada \\
nishal.silva@hotmail.com \\
+1 613 200 2030
\end{flushright}

\noindent Dear Hiring Manager,

\vspace{0.5em}
\noindent I am writing to apply for the Audio Software Engineer role on EA SPORTS UFC. My PhD research at the University of Trento was spent writing C++ audio systems against hard real-time budgets: I designed and shipped an audio pattern-detection engine that runs in 14 ms on a Raspberry Pi 4, holding F1 0.76 at 31.4\% CPU - 74x faster than the 1038 ms DTW baseline it replaced - published in the Journal of the Audio Engineering Society. I maintain Nebula, an open-source C++17 audio-feature-extraction library with per-feature ARM/Raspberry Pi latency benchmarks, and I build on JUCE, Elk Audio OS, VST, OSC and MIDI day to day. That combination - low-latency C++ audio code, DSP fundamentals, and a habit of measuring every millisecond - is exactly the discipline a live combat title like UFC demands: impacts, crowd reaction and dynamic mixing all have to land in real time, every time, with no room for a dropped frame of audio.

\vspace{0.5em}
\noindent Beyond the embedded work, I have shipped real-time audio under the least forgiving test there is: a live audience. As a postdoc I built a synchronization pipeline fusing audio, BCI/EEG and mixed-reality head/hand tracking across four musicians at once, then co-designed and ran two live multisensory concerts - reactive stage lighting, smoke and haptic feedback across ten Meta Quest 3 headsets - in front of a combined audience and performer group of 26 people, twice, without a system failure. That is the same pressure a shipped game system faces: audio has to hold up the instant it is in front of players, not just in a lab. I bring a musician's ear to that work as well, with 15+ years as a performing and session guitarist and a Trinity College London Rock \& Pop Guitar Grade 8, so audio quality is something I evaluate as a player, not only as an engineer.

\vspace{0.5em}
\noindent I have not worked inside a game studio before. This application is a deliberate pivot: I have spent a decade building real-time audio and sensor systems for research and performance, and I want to put that same rigor into a title played by millions rather than a lab audience of dozens. I am confident the gap is one of context, not capability - the C++, the DSP, and the real-time discipline transfer directly, and I am ready to learn the production pipeline and tooling specific to game audio quickly.

\vspace{0.5em}
\noindent I am a Canadian Permanent Resident and do not require sponsorship to work in Canada. I am based in Toronto, open to relocating to Vancouver, and available to start immediately. I would welcome the opportunity to discuss how my real-time audio engineering background can contribute to EA SPORTS UFC.

\vspace{1em}
\noindent Sincerely, \\
Nishal Stanislaus Silva

\end{document}
